%=============================================================================!
% 14.2  THE VIRIAL                                                            !
%=============================================================================!
\section{The virial}
\label{sec:pr-virial}

Atoms that pull on and push one another deliver momentum to the walls
differently from atoms that only bounce. This section finds the pressure
of interacting atoms from the equipartition theorem applied to their
positions, as Chapter~12 promised, and checks it against the push
measured in \cref{sec:pr-wall}.

\subsection*{Equipartition for positions}

\Cref{eq:sm-equipartition} says that $\ens{x\,\partial H/\partial x} =
\kB T$ for any coordinate $x$ of a system in the canonical ensemble whose
range is closed by a potential that rises without limit, walls included.
For atoms in a box, $H = K + U + U_{\mathrm w}$, with $U$ the energy of
their interactions and $U_{\mathrm w}$ that of the walls. $K$ depends on
the momenta alone, so the slope of $H$ along the coordinate $x_i$ of atom
$i$ is minus the force along it, $\partial H/\partial x_i = -F_{ix} -
F^{\mathrm w}_{ix}$, the first from the other atoms and the second from
the walls. The theorem for $x_i$ then reads $\ens{x_iF_{ix}} +
\ens{x_iF^{\mathrm w}_{ix}} = -\kB T$, and likewise for $y_i$ and $z_i$;
since $x_iF_{ix} + y_iF_{iy} + z_iF_{iz} = \vb{r}_i\cdot\vb{F}_i$, adding
the $3N$ of them gives
\begin{equation*}
   \Bigl\langle\sum_i\vb{r}_i\cdot\vb{F}_i\Bigr\rangle
   + \Bigl\langle\sum_i\vb{r}_i\cdot\vb{F}^{\mathrm w}_i\Bigr\rangle
   = -3N\kB T .
\end{equation*}
The first sum is the virial $\mathcal{W}$ of \cref{sec:pe-forces}, which
for atoms that interact in pairs is $\mathcal{W} = -\sum_{i<j}r_{ij}
\varphi'(r_{ij})$, \cref{eq:pe-virial}: positive when the atoms repel one
another, negative when they attract.

\subsection*{The walls' share}

The second sum comes only from atoms touching a wall. A wall pushes only
at right angles to itself, so an atom touching the wall at $x = L$ has
$\vb{r}_i\cdot\vb{F}^{\mathrm w}_i = x_iF^{\mathrm w}_{ix}$, with $x_i
\approx L$ and the force pointing back into the box; the sum over those
atoms is
\begin{equation*}
   \sum_ix_iF^{\mathrm w}_{ix} \approx L\sum_iF^{\mathrm w}_{ix} =
   -LF_\perp ,
\end{equation*}
with $F_\perp$ the force the atoms exert on the wall, equal and opposite
to the wall's force on them by the third law. On average $F_\perp$ is
$PA$, so this wall contributes $-PAL = -PV$. At the wall $x = 0$ the
atoms have $x_i \approx 0$ and contribute nothing; with the origin at the
centre of the box the two walls share $-PV$ equally instead
(\cref{ex:pr-centre}). The walls normal to $y$ and $z$ do the same, so the
walls' share is $-3PV$. That holds without error for hard walls. An atom
that has gone the depth $\delta$ into a soft wall adds a further
$-k\delta^2$, at $x = 0$ as at $x = L$, so soft walls enlarge the share by
a fraction of about $2\delta/L$. Putting it into the sum,
\begin{equation*}
   \ens{\mathcal{W}} - 3PV = -3N\kB T ,
\end{equation*}
and so
\begin{equation}
   PV = N\kB T + \tfrac13\ens{\mathcal{W}} .
   \label{eq:pr-virial}
\end{equation}
The first term is the ideal gas's, \cref{eq:pr-ideal}; the second is the
interactions' share, raising the pressure when the atoms repel one
another and lowering it when they attract. Since $\ens{2K} = 3N\kB T$ by
equipartition for the momenta, the instantaneous value
\begin{equation*}
   P = \frac{2K + \mathcal{W}}{3V}
\end{equation*}
averages to the pressure, and can be computed at every step from the
positions, velocities and forces the run already has.

\subsection*{Checked}

For the walled liquid of \cref{fig:pr-wall}(b), $(2K + \mathcal{W})/3V$
averages \SI{0.0628}{\giga\pascal} over the run (0.0621 and
\SI{0.0635}{\giga\pascal} in the halves), against \SI{0.0624}{\giga\pascal}
from the push on the walls: the two agree within their scatter. Of the
\SI{0.0628}{\giga\pascal}, \SI{0.0374}{\giga\pascal} is $N\kB T/V$ and
\SI{0.0254}{\giga\pascal} is $\ens{\mathcal{W}}/3V$: at this density the
repulsion between close neighbours outweighs the pull of distant ones.

\begin{keyidea}
Equipartition applied to the positions, $\ens{x\,\partial H/\partial x}
= \kB T$, gives $PV = N\kB T + \frac13\ens{\mathcal{W}}$: the ideal gas's
pressure plus the virial's share, positive for repulsion and negative for
attraction. The instantaneous pressure $(2K + \mathcal{W})/3V$ averages
to it and needs only the positions, velocities and forces.
\end{keyidea}

\begin{selfcheck}
Two argon atoms alone in a box sit \SI{3.0}{\angstrom} apart, closer than
the minimum of the Lennard-Jones potential at $2^{1/6}\sigma =
\SI{3.82}{\angstrom}$. Does their pair raise or lower the pressure from
that of two atoms that do not interact, and why?
\end{selfcheck}
